\begin{mistake}{2026-04-12}{35}

\begin{problemcontent}
设二元函数 \( u(x,y) \) 在点 \( M_0(x_0,y_0) \) 处取得极小值，且 \( u_{xx}(M_0) \)、\( u_{yy}(M_0) \) 存在。下列必成立的是：
(A) \( u_{xx}(M_0)\ge0,\ u_{yy}(M_0)\ge0 \)
(B) \( u_{xx}(M_0)>0,\ u_{yy}(M_0)>0 \)
\end{problemcontent}

\begin{solutioncontent}
选项 (A) 必成立。

**证明/分析**：
设 \( f(x) = u(x, y_0) \)。由于 \( u(x,y) \) 在 \( (x_0,y_0) \) 取极小值，则一元函数 \( f(x) \) 在 \( x_0 \) 也必取极小值。
根据一元函数极值的必要条件，若 \( f''(x_0) \) 存在，则必有 \( f''(x_0) \ge 0 \)。
而 \( f''(x_0) = u_{xx}(x_0, y_0) \)，故 \( u_{xx}(M_0) \ge 0 \)。同理可得 \( u_{yy}(M_0) \ge 0 \)。

**反例说明 (B)**：
对于 \( u(x,y) = x^4 + y^4 \)，在 \( (0,0) \) 取极小值，但 \( u_{xx}(0,0)=0, u_{yy}(0,0)=0 \)，不满足严格大于 0。
\end{solutioncontent}

\mistakeknowledge{
\begin{itemize}
  \item \textbf{核心原理}：多元极值的必要条件。若偏导存在，则梯度为零（一阶）；若二阶偏导存在，则海森矩阵（Hessian Matrix）半正定。
  \item \textbf{易错点}：混淆“必要条件”与“充分条件”。\( AC-B^2>0, A>0 \) 是极小值的充分条件，而非必要。
  \item \textbf{关联知识}：费马引理在高维空间的推广。
\end{itemize}}

\end{mistake}
